% !TeX root = ../technical_report.tex
\chapter{ONNX Runtime Acceleration and Quantisation Engine}
\label{ch:onnx_runtime_optimization}

In order to achieve stringent performance invariants and latency service level agreements within the production environment, the TaxonVision-MLOps deployment architecture transitions from high-level PyTorch computational graphs to deeply optimised, static ONNX (Open Neural Network Exchange) representations. The deployment strategy mandates zero memory allocation upon the hot inference path, dynamic eight-bit integer quantisation for memory compression, and deterministic execution providers that fail softly across hardware backends. This chapter details the comprehensive technical implementation of the ONNX exporter, the runtime session orchestration, the quantisation mechanics, and the visual interpretability paradigms encapsulated within the inference service pipeline.

\providecolor{TaxonPrimary}{HTML}{1976D2}
\providecolor{TaxonSecondary}{HTML}{388E3C}
\providecolor{TaxonAccent}{HTML}{F57C00}
\providecolor{TaxonDark}{HTML}{263238}

\begin{figure}[htbp]
    \centering
    \resizebox{\textwidth}{!}{%
    \begin{tikzpicture}[
        node distance=0.8cm and 0.5cm,
        font=\sffamily\footnotesize,
        % Core styles
        layer_box/.style={
            rectangle, rounded corners, draw=TaxonDark, fill=gray!10, thick,
            inner sep=12pt, align=center,
        },
        component_box/.style={
            rectangle, rounded corners, draw=#1, fill=#1!10, thick,
            minimum height=0.8cm, align=center,
            drop shadow={opacity=0.2, shadow xshift=1mm, shadow yshift=-1mm}
        },
        data_flow/.style={
            ->, >=Stealth, thick, TaxonDark
        },
        dashed_flow/.style={
            ->, >=Stealth, thick, dashed, TaxonDark
        },
        % Group backgrounds
        group_bg/.style={
            rectangle, rounded corners, draw=#1!50, dashed, fill=#1!5, inner sep=15pt, yshift=-5pt
        }
    ]

        % ==========================================
        % 1. PyTorch to ONNX Export Layer
        % ==========================================
        \node[component_box=TaxonPrimary] (pytorch_ckpt) {
            \textbf{PyTorch Checkpoint}\\
            \texttt{head.pt} + Frozen Backbone\\
            (FP32: $\sim$340 MB)
        };

        \node[component_box=TaxonPrimary, below=of pytorch_ckpt] (onnx_export) {
            \textbf{\texttt{torch.onnx.export}}\\
            Dynamic Axes: \texttt{batch\_size}\\
            Opset 17
        };

        \node[component_box=TaxonPrimary, below=of onnx_export] (graph_opt) {
            \textbf{Graph Optimisation}\\
            Constant Folding \&\\
            Operator Fusion (Conv+BN+ReLU)
        };

        \begin{scope}[on background layer]
            \node[group_bg=TaxonPrimary, fit=(pytorch_ckpt) (onnx_export) (graph_opt), 
                  label={[font=\sffamily\bfseries, text=TaxonPrimary]above:1. Computational Graph Export}] (export_group) {};
        \end{scope}

        \draw[data_flow] (pytorch_ckpt) -- (onnx_export);
        \draw[data_flow] (onnx_export) -- (graph_opt);

        % ==========================================
        % 2. INT8 Quantisation Engine
        % ==========================================
        \node[component_box=TaxonSecondary, right=2cm of pytorch_ckpt] (calib_data) {
            \textbf{Calibration Stream}\\
            Validation Observations
        };

        \node[component_box=TaxonSecondary, right=2cm of onnx_export] (quant_op) {
            \textbf{Static INT8 Calibration}\\
            Insert QDQ Nodes\\
            (QuantizeLinear / DequantizeLinear)
        };

        \node[component_box=TaxonSecondary, right=2cm of graph_opt] (quant_math) {
            \textbf{Precision Transformation}\\
            $q = \text{round}(x / S) + Z$\\
            (INT8: $\sim$86 MB, 4x Compression)
        };

        \begin{scope}[on background layer]
            \node[group_bg=TaxonSecondary, fit=(calib_data) (quant_op) (quant_math), 
                  label={[font=\sffamily\bfseries, text=TaxonSecondary]above:2. Quantisation Architecture}] (quant_group) {};
        \end{scope}

        \draw[data_flow] (calib_data) -- (quant_op);
        \draw[data_flow] (quant_op) -- (quant_math);
        \draw[data_flow] (graph_opt) -- (quant_op);

        % ==========================================
        % 3. Zero-Allocation Runtime Engine
        % ==========================================
        \node[component_box=TaxonDark, right=2cm of calib_data] (ort_session) {
            \textbf{\texttt{InferenceSession}}\\
            \texttt{CUDAExecutionProvider}\\
            Fail-soft to \texttt{CPUExecutionProvider}
        };

        \node[component_box=TaxonDark, below=of ort_session] (io_binding) {
            \textbf{Zero-Allocation Hot-Path}\\
            Pre-allocated \texttt{IOBinding} Buffers\\
            No Python GIL lock
        };

        \node[component_box=TaxonDark, below=of io_binding] (thread_pool) {
            \textbf{Concurrency}\\
            Threadpool Execution
        };

        \begin{scope}[on background layer]
            \node[group_bg=TaxonDark, fit=(ort_session) (io_binding) (thread_pool), 
                  label={[font=\sffamily\bfseries, text=TaxonDark]above:3. Inference Engine}] (runtime_group) {};
        \end{scope}

        \draw[data_flow] (ort_session) -- (io_binding);
        \draw[data_flow] (io_binding) -- (thread_pool);
        \draw[data_flow] (quant_math) -- (ort_session);

        % ==========================================
        % 4. EigenCAM Interpretability
        % ==========================================
        \node[component_box=TaxonAccent, below=1.5cm of quant_math] (cam_extract) {
            \textbf{Feature Map Extraction}\\
            Last Convolutional / Attention Layer
        };

        \node[component_box=TaxonAccent, right=of cam_extract] (cam_svd) {
            \textbf{Principal Component Saliency}\\
            Projection onto $v_1$ of Activation Matrix
        };

        \node[component_box=TaxonAccent, right=of cam_svd] (cam_visual) {
            \textbf{Visual Heatmap}\\
            Aligned with Taxonomic Prediction
        };

        \begin{scope}[on background layer]
            \node[group_bg=TaxonAccent, fit=(cam_extract) (cam_svd) (cam_visual), 
                  label={[font=\sffamily\bfseries, text=TaxonAccent]above:4. Forward-Only Interpretability (EigenCAM)}] (cam_group) {};
        \end{scope}

        \draw[data_flow] (cam_extract) -- (cam_svd);
        \draw[data_flow] (cam_svd) -- (cam_visual);
        
        % Connect Inference Engine to CAM
        \draw[dashed_flow] (thread_pool) -- (cam_visual) node[midway, right] {Prediction Output};
        \draw[dashed_flow] (runtime_group.south) |- (cam_extract.west) node[near start, left, align=right] {Activation\\Tensors};


    \end{tikzpicture}%
    }
    \caption{End-to-end architecture of the TaxonVision ONNX inference pipeline, illustrating the progression from PyTorch graph export and INT8 precision quantisation, through the zero-allocation runtime session, culminating in EigenCAM visual interpretability.}
    \label{fig:onnx_quantization_and_runtime}
\end{figure}


\section{Computational Graph Export Architecture}
\label{sec:onnx_graph_export}

In production workflows, the machine learning model must be strictly decoupled from the specific deep learning framework utilised during training in order to maximise deployment flexibility. ONNX serves as the de facto open-standard interoperability layer, enabling the model to be deployed using high-performance inference engines across diverse hardware targets without requiring the original PyTorch research and development environment. The module \texttt{src/taxon\_vision/inference/onnx\_exporter.py} encapsulates this translation through the \texttt{export\_to\_onnx} function.

The export process initiates by locking the neural network into evaluation mode, ensuring that stochastic operations such as dropout are disabled and batch normalisation statistics are frozen. Subsequently, the architecture employs the \texttt{torch.onnx.export} utility to execute a computational graph tracing procedure. By passing a dummy input tensor of shape \texttt{(1, 3, 224, 224)}, the PyTorch Just-In-Time (JIT) compiler records the exact sequence of mathematical operations executed upon the tensor, yielding a serialised static graph. Tracing is fundamentally preferred over scripting for this vision architecture, as it constructs a frozen execution graph that is highly amenable to low-level hardware optimisation, despite its inability to capture data-dependent control flow, which is notably absent in the standard feed-forward classification pipeline.

To balance advanced feature availability with broad hardware compatibility, the exporter specifies Operator Set (Opset) version 17. This modern opset provides robust support for the complex tensor manipulations inherent to transformer-based architectures and hierarchical vision backbones. Furthermore, dynamic axes are explicitly configured for the batch dimension:
\begin{minted}{python}
dynamic_axes={"input": {0: "batch_size"}, "logits": {0: "batch_size"}}
\end{minted}
This configuration preserves the capacity to process variable-sized workloads during inference. By defining the zero-th axis as dynamic, the compiled ONNX graph gracefully accommodates both single-sample real-time HTTP requests and high-throughput batch ingestion streams without necessitating computationally expensive graph recompilation.

In addition to structural mapping, the exporter enforces rigorous graph optimisation passes. The parameter \texttt{do\_constant\_folding=True} instructs the compiler to pre-calculate and collapse redundant deterministic operations -- such as scaling factors and bias additions -- into singular nodes. This constant folding mechanism, combined with redundant node elimination and kernel fusion, systematically minimises the computational depth of the final representation, directly translating into reduced memory bandwidth utilisation and accelerated inference latency. This graph manipulation fundamentally shifts computational complexity away from the inference endpoint and back towards the singular compilation phase, thereby guaranteeing that the live server expends absolute minimum cycles parsing neural architectures. The resulting binary artifact encapsulates not just the model weights, but a fully deterministic execution plan optimised specifically for rapid tensor processing.

\section{Zero-Allocation Inference Engine and Session Management}
\label{sec:zero_allocation_inference_engine}

The operational constraints of the biological observation ingestion pipeline mandate strict sub-25 millisecond 95th-percentile (p95) latency bounds. To satisfy this Service Level Objective (SLO), the inference engine, detailed within \texttt{src/taxon\_vision/service/inference.py}, enforces a rigorous zero-allocation invariant upon the hot execution path. Dynamic heap allocations trigger non-deterministic garbage collection cycles, which introduce unacceptable latency spikes during high-throughput load conditions. By eliminating dynamic heap allocations on the hot path, the system mathematically guarantees predictable response times regardless of query burst density.

To eliminate per-request instantiation overhead, the service layer heavily leverages process-level memoisation via the \texttt{@lru\_cache(maxsize=1)} decorator. Critical objects, including the \texttt{TaxonClassifier}, \texttt{ConformalPredictionEngine}, \texttt{EnergyOODDetector}, and the authoritative taxa catalog, are instantiated precisely once during application startup. Crucially, the image preprocessing pipeline is also statically cached. The \texttt{get\_image\_transform} routine bundles \texttt{Resize}, \texttt{ToTensor}, and \texttt{Normalize} operations into a single immutable \texttt{torchvision.transforms.Compose} instance. By avoiding the recurrent construction of these transformation objects, the system bypasses severe memory fragmentation and continuous allocation overhead.

When a raw byte payload arrives via the network interface, it is wrapped in an in-memory \texttt{BytesIO} stream, decoded as a PIL image, transformed using the pre-cached pipeline, and augmented with a batch dimension. This methodology ensures that the hot path operates entirely upon pre-allocated buffers and zero-gradient computational spaces. No intermediate tensors are permitted to escape the local functional scope, ensuring pristine memory discipline.

Within the broader ONNX Runtime session management context, the architecture implements sophisticated Execution Provider (EP) negotiation. The inference session is configured to prioritise the \texttt{CUDAExecutionProvider}, securing pinned GPU memory for maximal throughput. However, to guarantee uninterrupted service during transient GPU time-slice preemption or hardware faults, the session configuration provides a seamless, fail-soft fallback mechanism to the \texttt{CPUExecutionProvider}. This multi-tier hardware targeting ensures resilient operation regardless of underlying node health. If the NVIDIA kernel driver crashes or reports an out-of-memory exception, the ONNX execution context transparently migrates the workload to system RAM, gracefully degrading throughput rather than hard-failing the user request.

Performance tuning extends deeply into the runtime threadpool configuration. The engine carefully balances \texttt{intra\_op\_num\_threads} and \texttt{inter\_op\_num\_threads}. Intra-operation parallelism accelerates individual mathematical kernels (such as large matrix multiplications) across multiple CPU cores, whereas inter-operation parallelism permits concurrent execution of independent graph branches. Moreover, the session mode is explicitly evaluated, contrasting \texttt{ORT\_SEQUENTIAL} evaluation -- which minimises synchronisation overhead for linear topologies -- against \texttt{ORT\_PARALLEL} evaluation for branched neural architectures. This comprehensive orchestration yields an inference engine capable of sustaining immense throughput whilst remaining entirely immune to memory leak phenomena. The precise orchestration of these threading paradigms directly prevents thread exhaustion under massive parallel ingress, fortifying the web server against cascade failures.

\section{INT8 Quantisation Engine and Precision Transformation}
\label{sec:quantization_and_cam}

The deployment of deep neural networks in resource-constrained or high-density containerised environments necessitates aggressive memory compression. The TaxonVision architecture implements an advanced eight-bit integer (INT8) quantisation engine to map 32-bit floating-point (FP32) network parameters into a compact discrete space. This mathematical transformation yields a four-fold reduction in Video Random Access Memory (VRAM) consumption and significantly accelerates matrix multiplication bandwidth. By shrinking the classification weights from 32-bit floats to 8-bit integers, the pipeline enables multiple models to reside simultaneously within the tight VRAM limits of edge appliances.

Quantisation mapping can be executed through both symmetric and asymmetric mathematical formulations. In the asymmetric approach, the floating-point tensor $x$ is mapped to a quantised integer representation $q$ using a scale factor $S$ and an integer zero-point $Z$:
\begin{equation}
    q = \text{round}\left(\frac{x}{S}\right) + Z
\end{equation}
The inclusion of the zero-point $Z$ ensures that the floating-point value $0.0$ maps exactly to a quantised integer, a crucial property for preserving the zero-sparsity introduced by activation functions such as the Rectified Linear Unit (ReLU). Alternatively, symmetric quantisation forces $Z=0$, simplifying the arithmetic pipeline at the cost of a slightly contracted dynamic range. For the classification weights, dynamic quantisation provides rapid compression by calculating scaling factors on-the-fly during inference, whilst static quantisation relies upon a \texttt{TaxonCalibrationDataReader} to measure the absolute distribution of activations offline, enabling the complete quantisation of both weights and activations. Static methodologies inherently produce faster inference graphs, as all threshold clipping logic is pre-computed and hardcoded into the compiled binary.

This precision transformation naturally introduces a trade-off between execution latency and classification accuracy. The rounding operations inherently discard low-magnitude weight precision, which can precipitate accuracy degradation, particularly amongst rare, long-tailed biological classes where decision boundaries are narrow. However, the system mitigates this risk through its robust training topology. The adoption of a Class-Balanced Loss function during training ensures that the structural representations of minority taxa are deeply ingrained into the feature space. Furthermore, the Split Conformal Prediction engine buffers minor logit perturbations by dynamically expanding the prediction set cardinality whenever the underlying quantised confidence deteriorates. Consequently, the pipeline achieves maximal acceleration while preserving strict epistemic safety guarantees, validating that quantitative performance metrics do not catastrophically undermine qualitative ecological intelligence.

\section{Visual Interpretability via Saliency Attribution}
\label{sec:visual_interpretability}

Providing diagnostic transparency to domain experts requires the generation of high-fidelity visual explainability artifacts. Class Activation Mapping (CAM) highlights the precise spatial regions within an organism photograph that most significantly influenced the model's taxonomic conclusion. As detailed within \texttt{src/taxon\_vision/inference/cam.py}, the architectural constraints of the high-throughput inference engine impose strict technical limitations upon the implementation of these attribution methods. A visually interpretable heatmap is absolutely mandatory for establishing trust with the scientific community, yet generating such graphics often imposes severe computational penalties on the deployment infrastructure.

Classical attribution algorithms, most notably Grad-CAM, are strictly prohibited within this production environment. Grad-CAM mandates a backward gradient pass through the entire neural backbone. This operation triples the baseline inference latency, overtly violating the sub-25 millisecond operational SLO. More critically, gradient backpropagation is entirely incompatible with quantised INT8 ONNX Runtime inference engines, wherein the gradient computation graphs are deliberately stripped out to minimise the binary footprint. The architecture therefore demands forward-only attribution mechanics that completely bypass automatic differentiation.

The strategic target for this forward-only interpretability is EigenCAM. EigenCAM generates class-discriminative visual saliency heatmaps by computing the principal components of the deepest feature map representations. By extracting the dominant eigenvectors from the final convolutional or attention layers, EigenCAM identifies the most structurally relevant features without necessitating any backward pass or parameter tuning. This property is exceptionally advantageous for vision transformers and modern hierarchical backbones, where traditional gradient-based CAM methods frequently suffer from severe gradient diffusion and numerical instability. Furthermore, EigenCAM requires precisely zero parameter tuning, enabling it to operate universally across diverse architectural families without necessitating laborious grid searches for optimal regularisation constants.

Prior to the finalisation of multi-head feature-map forward hooking within the ONNX engine, the system implements an interim radial saliency placeholder to fulfill the immediate HTMX/Jinja2 User Interface contract. This highly optimised surrogate executes in under two milliseconds. It functions by extracting the spatial dimensions of the input image and constructing an open mesh grid centered upon the geometric centroid. It calculates the normalised radial Euclidean distances and applies an exponential Gaussian attenuation function:
\begin{equation}
    \text{Heatmap}(x, y) = \exp\left( \frac{-\sqrt{(x - c_x)^2 + (y - c_y)^2}}{w / 3.0} \right)
\end{equation}
Finally, the resulting activation surface is subjected to min-max scaling, translating the raw values into the strictly bounded $[0.0, 1.0]$ interval. This yields a clean, \texttt{float32} spatial mask that mimics structural saliency, preserving absolute latency invariants while the engineering foundation is laid for the rigorous EigenCAM integration.
